\label{sec:mechanisms_body}

\subsection{Decomposition by question type}\label{ssec:mech_by_type}\label{ssec:mech_qtype_anomaly}

Estimating the AI$\times$Post interaction separately for each
question type (tag and week fixed effects) orders the
coefficients as the substitutability theory predicts:
short-code ($-0.31$), how-to ($-0.28$), long-code ($-0.21$),
debugging-simple ($-0.14$), other-conceptual ($-0.08$,
marginal), advanced-architecture ($-0.16$, $p=0.003$), and
version-environment-specific (null, $-0.03$, $p=0.54$). The
version-environment-specific null---the genuinely clean
control---is the only category with no detected decline. That
the declared boundary category advanced-architecture shows a
sizeable effect confirms its a-priori boundary status
(\S\ref{ssec:data_ai_answerability}); because the binary
classifier codes it $\text{Sub}_k=0$, it contributes a negative
shift to the control group and mechanically attenuates the full
sample, which is why excluding it strengthens the funnel DDD
from $-0.108$ to $-0.140$ (full table in the replication
package). Replacing the binary indicator with the continuous
embedding score gives $-0.039$ ($p=0.004$), about a third the
binary magnitude as expected from shrinking the contrast; the
qualitative conclusion is unchanged.

\subsection{Composition versus within-cell
quality}\label{ssec:mech_complementarity}\label{ssec:mech_resolvability}

A complementarity reading holds that users keep posting when an
LLM answer is unsatisfactory, so the \emph{surviving} questions
should be more complex. Descriptively this holds: within tags,
post-period body length rises ($+303$ characters, $p<10^{-15}$)
and accepted-answer share falls ($-0.065$, $p<10^{-15}$). But
the Post$\times$AI-answerability interactions are null
($p$ between $0.32$ and $0.79$), so the compositional shift is
\emph{general across tags} rather than concentrated where
answerability is highest---it is descriptive selection, not a
causally estimated quality effect along the AI dimension.

Separating selection from a within-cell quality channel is the
contribution here. \citet{xue2026chatgpt} report higher average
question quality post-ChatGPT under a cross-platform design that
admits both channels. Re-fitting the DDD on six resolvability
and quality outcomes under fine-grained (tag$\times$question-type)
and week fixed effects---which absorb the composition shift---five
of six estimates are null (accepted share $-0.003$, $p=0.45$;
answered share $-0.001$, $p=0.85$; analogous nulls elsewhere);
the only significant coefficient is a small \emph{rise} in
closed share ($+0.005$, $p<0.001$). The descriptively documented
quality improvement is therefore consistent with selection into
the survivor population (the easier posts displaced into private
channels) rather than with within-cell upgrading. This is not a
contradiction of \citet{xue2026chatgpt}: the cross-platform ATT
(aggregate declines of $\sim$15--25\%;
\citealp{delriochanona2024large}) and the within-platform
$\sim$5\% slice estimated here are complementary objects, not
competing estimates of the same quantity (full mapping and
detail tables in the replication package).

\subsection{Where in the workflow does substitution
operate?}\label{ssec:mech_copilot}

The June 2022 GitHub Copilot placebo is positive ($+0.058$,
$p=0.011$): the Stack Overflow Q\&A margin responded to ChatGPT
but not to Copilot. Two non-exclusive readings are consistent
with the data and we cannot separate them: Copilot operated
inside the IDE at the code-writing moment rather than the
search-for-help moment that ends in a public query, and its 2022
user base was much smaller than ChatGPT's. The \emph{positive}
sign is the informative part (a shared secular trend cannot
produce it), but its magnitude rests on the smaller 2022
population, so the test has limited power; we read it as
supporting evidence on the decision-point reading, not a
stand-alone mechanism, and present any interface-location
interpretation as a post-hoc hypothesis the design does not
identify.
